%!TeX root=../gardentop.tex
\chapter{The Cry in the Corridor} 
	
\begin{figure}[t!]
\centering
\includegraphics[width=\linewidth]{may}
\end{figure}

 \lettrine[lines=6]{A}{t} first each day which passed by for Mary Lennox was exactly like the others. Every morning she awoke in her tapestried room and found Martha kneeling upon the hearth building her fire; every morning she ate her breakfast in the nursery which had nothing amusing in it; and after each breakfast she gazed out of the window across to the huge moor which seemed to spread out on all sides and climb up to the sky, and after she had stared for a while she realized that if she did not go out she would have to stay in and do nothing—and so she went out. She did not know that this was the best thing she could have done, and she did not know that, when she began to walk quickly or even run along the paths and down the avenue, she was stirring her slow blood and making herself stronger by fighting with the wind which swept down from the moor. She ran only to make herself warm, and she hated the wind which rushed at her face and roared and held her back as if it were some giant she could not see. But the big breaths of rough fresh air blown over the heather filled her lungs with something which was good for her whole thin body and whipped some red colour into her cheeks and brightened her dull eyes when she did not know anything about it.

But after a few days spent almost entirely out of doors she wakened one morning knowing what it was to be hungry, and when she sat down to her breakfast she did not glance disdainfully at her porridge and push it away, but took up her spoon and began to eat it and went on eating it until her bowl was empty.

»Tha' got on well enough with that this mornin', didn't tha'?« said Martha.

»It tastes nice to-day,« said Mary, feeling a little surprised herself.

»It's th' air of th' moor that's givin' thee stomach for tha' victuals,« answered Martha. »It's lucky for thee that tha's got victuals as well as appetite. There's been twelve in our cottage as had th' stomach an' nothin' to put in it. You go on playin' you out o' doors every day an' you'll get some flesh on your bones an' you won't be so yeller.«

»I don't play,« said Mary. »I have nothing to play with.«

»Nothin' to play with!« exclaimed Martha. »Our children plays with sticks and stones. They just runs about an' shouts an' looks at things.«

Mary did not shout, but she looked at things. There was nothing else to do. She walked round and round the gardens and wandered about the paths in the park. Sometimes she looked for Ben Weatherstaff, but though several times she saw him at work he was too busy to look at her or was too surly. Once when she was walking toward him he picked up his spade and turned away as if he did it on purpose.

One place she went to oftener than to any other. It was the long walk outside the gardens with the walls round them. There were bare flower-beds on either side of it and against the walls ivy grew thickly. There was one part of the wall where the creeping dark green leaves were more bushy than elsewhere. It seemed as if for a long time that part had been neglected. The rest of it had been clipped and made to look neat, but at this lower end of the walk it had not been trimmed at all.

A few days after she had talked to Ben Weatherstaff Mary stopped to notice this and wondered why it was so. She had just paused and was looking up at a long spray of ivy swinging in the wind when she saw a gleam of scarlet and heard a brilliant chirp, and there, on the top of the wall, perched Ben Weatherstaff's robin redbreast, tilting forward to look at her with his small head on one side.

»Oh!« she cried out, »is it you—is it you?« And it did not seem at all queer to her that she spoke to him as if she was sure that he would understand and answer her.

He did answer. He twittered and chirped and hopped along the wall as if he were telling her all sorts of things. It seemed to Mistress Mary as if she understood him, too, though he was not speaking in words. It was as if he said:

»Good morning! Isn't the wind nice? Isn't the sun nice? Isn't everything nice? Let us both chirp and hop and twitter. Come on! Come on!«

Mary began to laugh, and as he hopped and took little flights along the wall she ran after him. Poor little thin, sallow, ugly Mary—she actually looked almost pretty for a moment.

»I like you! I like you!« she cried out, pattering down the walk; and she chirped and tried to whistle, which last she did not know how to do in the least. But the robin seemed to be quite satisfied and chirped and whistled back at her. At last he spread his wings and made a darting flight to the top of a tree, where he perched and sang loudly.

That reminded Mary of the first time she had seen him. He had been swinging on a tree-top then and she had been standing in the orchard. Now she was on the other side of the orchard and standing in the path outside a wall—much lower down—and there was the same tree inside.

»It's in the garden no one can go into,« she said to herself. »It's the garden without a door. He lives in there. How I wish I could see what it is like!«

She ran up the walk to the green door she had entered the first morning. Then she ran down the path through the other door and then into the orchard, and when she stood and looked up there was the tree on the other side of the wall, and there was the robin just finishing his song and beginning to preen his feathers with his beak.

»It is the garden,« she said. »I am sure it is.«

She walked round and looked closely at that side of the orchard wall, but she only found what she had found before—that there was no door in it. Then she ran through the kitchen-gardens again and out into the walk outside the long ivy-covered wall, and she walked to the end of it and looked at it, but there was no door; and then she walked to the other end, looking again, but there was no door.

»It's very queer,« she said. »Ben Weatherstaff said there was no door and there is no door. But there must have been one ten years ago, because Mr~Craven buried the key.«

This gave her so much to think of that she began to be quite interested and feel that she was not sorry that she had come to Misselthwaite Manor. In India she had always felt hot and too languid to care much about anything. The fact was that the fresh wind from the moor had begun to blow the cobwebs out of her young brain and to waken her up a little.

She stayed out of doors nearly all day, and when she sat down to her supper at night she felt hungry and drowsy and comfortable. She did not feel cross when Martha chattered away. She felt as if she rather liked to hear her, and at last she thought she would ask her a question. She asked it after she had finished her supper and had sat down on the hearth-rug before the fire.

»Why did Mr~Craven hate the garden?« she said.

She had made Martha stay with her and Martha had not objected at all. She was very young, and used to a crowded cottage full of brothers and sisters, and she found it dull in the great servants' hall down-stairs where the footman and upper-housemaids made fun of her Yorkshire speech and looked upon her as a common little thing, and sat and whispered among themselves. Martha liked to talk, and the strange child who had lived in India, and been waited upon by »blacks,« was novelty enough to attract her.

She sat down on the hearth herself without waiting to be asked.

»Art tha' thinkin' about that garden yet?« she said. »I knew tha' would. That was just the way with me when I first heard about it.«

»Why did he hate it?« Mary persisted.

Martha tucked her feet under her and made herself quite comfortable.

»Listen to th' wind wutherin' round the house,« she said. »You could bare stand up on the moor if you was out on it to-night.«

Mary did not know what »wutherin'« meant until she listened, and then she understood. It must mean that hollow shuddering sort of roar which rushed round and round the house as if the giant no one could see were buffeting it and beating at the walls and windows to try to break in. But one knew he could not get in, and somehow it made one feel very safe and warm inside a room with a red coal fire.

»But why did he hate it so?« she asked, after she had listened. She intended to know if Martha did.

Then Martha gave up her store of knowledge.

»Mind,« she said, »Mrs~Medlock said it's not to be talked about. There's lots o' things in this place that's not to be talked over. That's Mr~Craven's orders. His troubles are none servants' business, he says. But for th' garden he wouldn't be like he is. It was Mrs~Craven's garden that she had made when first they were married an' she just loved it, an' they used to 'tend the flowers themselves. An' none o' th' gardeners was ever let to go in. Him an' her used to go in an' shut th' door an' stay there hours an' hours, readin' an' talkin'. An' she was just a bit of a girl an' there was an old tree with a branch bent like a seat on it. An' she made roses grow over it an' she used to sit there. But one day when she was sittin' there th' branch broke an' she fell on th' ground an' was hurt so bad that next day she died. Th' doctors thought he'd go out o' his mind an' die, too. That's why he hates it. No one's never gone in since, an' he won't let any one talk about it.«

Mary did not ask any more questions. She looked at the red fire and listened to the wind »wutherin'.« It seemed to be »wutherin'« louder than ever.

At that moment a very good thing was happening to her. Four good things had happened to her, in fact, since she came to Misselthwaite Manor. She had felt as if she had understood a robin and that he had understood her; she had run in the wind until her blood had grown warm; she had been healthily hungry for the first time in her life; and she had found out what it was to be sorry for some one. She was getting on.

But as she was listening to the wind she began to listen to something else. She did not know what it was, because at first she could scarcely distinguish it from the wind itself. It was a curious sound—it seemed almost as if a child were crying somewhere. Sometimes the wind sounded rather like a child crying, but presently Mistress Mary felt quite sure that this sound was inside the house, not outside it. It was far away, but it was inside. She turned round and looked at Martha.

»Do you hear any one crying?« she said.

Martha suddenly looked confused.

»No,« she answered. »It's th' wind. Sometimes it sounds like as if some one was lost on th' moor an' wailin'. It's got all sorts o' sounds.«

»But listen,« said Mary. »It's in the house—down one of those long corridors.«

And at that very moment a door must have been opened somewhere down-stairs; for a great rushing draft blew along the passage and the door of the room they sat in was blown open with a crash, and as they both jumped to their feet the light was blown out and the crying sound was swept down the far corridor so that it was to be heard more plainly than ever.

»There!« said Mary. »I told you so! It is some one crying—and it isn't a grown-up person.«

Martha ran and shut the door and turned the key, but before she did it they both heard the sound of a door in some far passage shutting with a bang, and then everything was quiet, for even the wind ceased »wutherin'« for a few moments.

»It was th' wind,« said Martha stubbornly. »An' if it wasn't, it was little Betty Butterworth, th' scullery-maid. She's had th' toothache all day.«

But something troubled and awkward in her manner made Mistress Mary stare very hard at her. She did not believe she was speaking the truth.